\section{Discussion and Conclusion}\label{sec:discussion}

This study presents a framework for text difficulty estimation and content adaptation using Large Language Models, addressing the fundamental challenge in language learning of matching content difficulty to learner proficiency. The framework integrates three components, each supported by specific evaluation metrics. The Pairwise Mismatch (PM) metric evaluates the preservation of difficulty ordering between text pairs, addressing the relative nature of text difficulty assessment. For content simplification, we developed a dual-metric approach combining simplification accuracy (A) and semantic similarity (S) through a weighted-score $(w\text{-Score} = 2(w_1A \times w_2S)/(w_1A + w_2S))$, with empirically determined weights $(w_1 = 0.80, w_2 = 0.20)$.

Our experiments across three French language datasets reveal several key insights about difficulty estimation. While larger models like GPT-4o achieved the highest absolute performance, specialized models like CamemBERT and Mistral-7B demonstrated surprisingly competitive results despite their significantly smaller size. This suggests that fine-tuned, specialized models can match or even outperform state-of-the-art models in specific tasks. The Pairwise Mismatch metric proved particularly valuable in highlighting these nuances, revealing cases where models achieved similar accuracy but demonstrated notably different ordering preservation capabilities. The effectiveness of system prompts across all models, including non-instruction-tuned ones like Davinci and Babbage, indicates that explicit task descriptions consistently enhance model performance, even for fine-tuned models.

The text simplification results revealed unexpected patterns in model behavior. The fine-tuned Mistral-7B model emerged as the most effective, balancing simplification accuracy with semantic preservation. Surprisingly, larger models like GPT-4o and GPT-4o-Mini showed less favorable performance, adopting more aggressive simplification strategies that often removed content rather than reformulating it. This contrasts with smaller models like GPT-3.5 and Mistral, which favored more conservative approaches using synonym substitution and syntactic restructuring. These findings carry important implications for environmental sustainability, as they demonstrate that effective text simplification can be achieved with smaller, more energy-efficient models.

The recommendation system evaluation demonstrates the advantages of integrating LLM-based embeddings into traditional collaborative filtering architectures. By combining difficulty estimation and content adaptation with personalized recommendations, the system addresses a critical aspect of extensive reading: maintaining learner engagement through appropriately challenging content. The performance gains were most pronounced in datasets with comprehensive textual descriptions, suggesting that the quality of textual content significantly influences the effectiveness of LLM-based embeddings. This integrated approach enables a more sophisticated matching of content to both learner proficiency and interests.

\textcolor{red}{While our empirical validation was conducted entirely on French reading corpora (and English/multilingual metadata in the recommendation benchmarks), the architectural formulation of our pipeline ($PM_{\text{norm}}$, harmonic $w\text{-Score}$, and LLM-initialized LightGCN propagation) does not rely on French-specific morphological heuristics and can be transferred to other CEFR-aligned languages by substituting the backbone encoder and parallel training pairs.} The approach provides a methodology for content adaptation that can be integrated into computer-assisted language learning systems, with empirical evidence supporting its effectiveness across different text genres and difficulty levels. The results particularly highlight the potential of smaller, specialized models in providing computationally efficient solutions without significantly compromising performance.

Several limitations require consideration for future work. The framework's performance depends on training data quality and diversity, with potential biases in both difficulty assessment and simplification processes. \textcolor{red}{While our double-blind human evaluation on 150 stratified sentence pairs (Section~\ref{subsubsec:human_eval}) validates sentence-level grammaticality, semantic preservation, and perceived simplification gain, larger-scale longitudinal classroom studies remain necessary to measure long-term vocabulary acquisition and reading comprehension gains.} Technical constraints include processing requirements and the need for consistent cross-language evaluation metrics. Future research directions include longitudinal studies examining the impact on language acquisition outcomes in extensive reading contexts, investigation of discourse-level features in difficulty assessment, and cross-linguistic validation studies to examine the framework's applicability across different language pairs. Additionally, exploring knowledge distillation techniques could further optimize the efficiency of smaller models while maintaining performance levels.